\documentclass[11pt,a4paper]{article}
\usepackage[T1]{fontenc}
\usepackage[utf8]{inputenc}
\usepackage{lmodern}
\usepackage[a4paper,margin=25mm,headheight=14pt]{geometry}
\usepackage{amsmath,amssymb,amsthm,mathtools}
\usepackage{microtype,booktabs,tabularx,array,enumitem}
\usepackage{fancyhdr,xurl,hyperref}
\hypersetup{hidelinks,pdftitle={Path-Sensitive Small Divisors in Hahn--Fuchsian Systems},pdfauthor={Research report prepared for Vladimir Reshetnikov}}
\pagestyle{fancy}
\fancyhf{}
\fancyhead[L]{\small Path-sensitive Hahn--Fuchsian normalization}
\fancyhead[R]{\small 29 September 2026}
\fancyfoot[C]{\thepage}
\setlength{\emergencystretch}{2em}
\setlength{\parskip}{3pt}
\setlist{itemsep=3pt,topsep=5pt}
\numberwithin{equation}{section}
\newtheorem{theorem}{Theorem}[section]
\newtheorem{proposition}[theorem]{Proposition}
\newtheorem{lemma}[theorem]{Lemma}
\newtheorem{corollary}[theorem]{Corollary}
\theoremstyle{definition}
\newtheorem{definition}[theorem]{Definition}
\newtheorem{example}[theorem]{Example}
\theoremstyle{remark}
\newtheorem{remark}[theorem]{Remark}
\newcommand{\R}{\mathbb R}
\newcommand{\C}{\mathbb C}
\newcommand{\N}{\mathbb N}
\newcommand{\K}{\mathcal K}
\newcommand{\A}{\mathcal A}
\newcommand{\G}{\mathcal G}
\newcommand{\supp}{\operatorname{supp}}
\newcommand{\diag}{\operatorname{diag}}
\newcommand{\Mat}{\operatorname{Mat}}
\newcommand{\lc}{\operatorname{lc}_{\mathrm{dp}}}
\newcommand{\Paths}{\operatorname{Paths}}
\newcommand{\abs}[1]{\left|#1\right|}
\newcommand{\norm}[1]{\left\|#1\right\|}
\newcommand{\dplog}[1]{L^{[#1]}}
\newcommand{\snap}{\texttt{9250bbf8af80dacf7b252d9d0320b122c80961ba}}
\newcommand{\repo}{\texttt{VladimirReshetnikov/ProveIt}}
\newcommand{\eps}{\varepsilon}

\title{\textbf{Path-Sensitive Small Divisors\\in Hahn--Fuchsian Systems}\\[5pt]
\large Exact convergence criteria, resonance amplification,\\and stable analytic renormalization}
\author{Research report prepared for Vladimir Reshetnikov}
\date{29 September 2026}

\begin{document}
\maketitle
\thispagestyle{plain}
\begin{abstract}
A recent ProveIt manuscript characterizes universal absolute convergence of
Hahn--Fuchsian normalizers when every exponent in a prescribed semigroup is
an allowed input. It leaves open the corresponding problem for a smaller
input support and restricted matrix directions. We resolve that problem for
strictly upper-triangular perturbations of a real diagonal matrix, with
independently variable coefficients on each permitted edge. The exact
criterion involves only exponent sums along actual directed matrix paths:
no endpoint spectral difference may be a left accumulation point of its
path sumset. Accumulations created by algebraically impossible products are
irrelevant. Exact resonances are allowed.

The proof gives an explicit polynomial-in-logarithm path kernel, an exact
weighted multilinear norm criterion, and a conversion to the usual
normalized log-free gauge. It also yields arbitrarily small,
bounded-support counterexamples whenever the criterion fails. Two
three-dimensional examples distinguish the path criterion from both a
whole-semigroup test and an edge-by-edge test. An arbitrary-rank resonant
chain exhibits a sharp loss of moments: coefficients $c_n$ require exactly
$\sum n^{r-1}|c_n|<\infty$ in dimension $r$, although the original input
only requires $\sum|c_n|<\infty$.

For that chain, a stable basepoint-normalized analytic solution exists under
the weaker input condition and realizes every fixed finite formal prefix.
We prove a uniform $N\asymp\log(1/x)$ crossover with an explicit error bound.
All specialized arguments are supplied. Exact finite checks and numerical
illustrations accompany the article; they are not a Lean formalization.
The research additions are relative to the inspected repository questions;
worldwide priority has not been established.
\end{abstract}

\newpage
\setcounter{tocdepth}{1}
\tableofcontents
\vspace{12pt}
\section*{Notation at a glance}
\begin{center}
\begin{tabularx}{\textwidth}{@{}lX@{}}
\toprule
Symbol & Meaning\\
\midrule
$D$, $L$ & Euler derivation $x\,d/dx$ and the formal logarithm, $DL=1$.\\
$\Lambda$, $\G$ & Real diagonal constant matrix and the acyclic edge pattern.\\
$E_p$, $\rho_p$ & Exponent sumset and spectral difference for a path $p$.\\
$T_\beta$, $Q_p$ & Normalized polynomial inverse and its composed path kernel.\\
$H$, $P$ & Canonical polynomial-log factor and normalized log-free gauge.\\
$M_p(w)$, $\eta_p$ & Weighted path norm and the nonresonant path gap.\\
\bottomrule
\end{tabularx}
\end{center}
\newpage

\section{The research question and the scope of the answer}
\label{sec:orientation}

\subsection{A specific question in the repository}
The ProveIt article \emph{Finite Resonance Certificates and Sharp Analytic
Normalization for Hahn--Fuchsian Systems} studies
\[
 xY'=(A_0+A_+(x))Y
\]
with positive, well-ordered real exponents. Its universal convergence theorem
quantifies over all coefficient series supported in an entire exponent
semigroup. Section 11 then asks whether that criterion remains necessary
when the input is restricted to a smaller generating support. It explicitly
suggests a triangular three-dimensional system with two permitted matrix
units and an accumulating two-letter sumset \cite{repo-hahn}.

This is not a cosmetic change of quantifiers. A semigroup remembers which
exponents can be added, but forgets whether the associated matrix products
are nonzero. For instance, the repeated product of a matrix unit $E_{12}$
is zero, whereas $E_{12}E_{23}=E_{13}$ is not. A spectral accumulation arising
only from the first kind of word need not obstruct normalization. Conversely,
a bad sum formed by the second kind can be invisible in both individual
coefficient series.

We solve this support-sensitive problem for every finite acyclic matrix
pattern, not only the suggested three-dimensional example. After a
permutation of vertices, such a pattern is strictly upper triangular.
The constant term is diagonal with real eigenvalues; coefficients on
different edges are independently variable. These are genuine restrictions,
and the conclusions below retain them explicitly.

\subsection{What is proved, and what is not}
The principal results are as follows.
\begin{center}
\begin{tabularx}{\textwidth}{@{}>{\raggedright\arraybackslash}p{.28\textwidth}X@{}}
\toprule
Question or issue & Result in this article\\
\midrule
Smaller input supports and allowed matrix directions &
An exact path--spectrum equivalence, including exact resonances;
Theorems~\ref{thm:gap} and \ref{thm:gauge-equivalence}.\\[3pt]
Weighted convergence when spectral gaps vanish &
An exact necessary and sufficient path-kernel criterion for independent
weighted $\ell^1$ input classes; Theorem~\ref{thm:weighted}.\\[3pt]
How severe can resonance amplification be? &
In dimension $r$, an explicit chain has the exact moment threshold
$n^{r-1}$; Theorem~\ref{thm:chain}.\\[3pt]
What replaces a divergent canonical expansion? &
A basepoint-normalized realization for the chain, with finite-prefix
asymptotics and a uniform crossover certificate;
Theorems~\ref{thm:analytic-chain} and \ref{thm:crossover}.\\
\bottomrule
\end{tabularx}
\end{center}

We do not settle the smaller-support problem for arbitrary nontriangular
matrices, fixed nonzero Jordan blocks, or coupled coefficient directions.
We do not construct a multiplicative summation operation on an unrestricted
Hahn field. The analytic realization theorem concerns the stated chain,
not every formally divergent transseries equation. The distinction between
universal independent inputs and correlated inputs is demonstrated
explicitly in Section~\ref{sec:cancellation}.

\subsection{Relation to established theory}
The Frobenius method, resonant Fuchsian normal forms, and the integration of
nilpotent normal forms are classical; see Ilyashenko--Yakovenko
\cite{iy,yakovenko}. Well-based series and their differential calculus belong
to the established theory of transseries \cite{vdh}. Neither that general
architecture nor the finite-dimensional inverse of
$\beta+\partial_L$ is claimed as a new theorem.

Gontsov--Goryuchkina give a convergence theorem for generalized power-series
solutions of polynomial ordinary differential equations; their displayed
class has $\operatorname{Re}s_k\to+\infty$ \cite{gg}. Our coefficient data can
instead have infinitely many exponents below a finite cut and are not
assumed polynomial in $x$. The small-divisor work of
Gontsov--Goryuchkina--Lastra concerns generalized solutions of
$q$-difference equations \cite{ggl}. It is relevant context, not a source of
the graph-specific equivalence proved here.

The proposed additions are the sharp fixed-pattern criterion, its exact
weighted version, the all-radii counterexample consequence, and the
resonant-chain moment and crossover package. The earlier repository
manuscript already supplies a whole-semigroup criterion, finite resonance
ancestry, and a rank-one accumulation model. We build on its questions
rather than relabel those results as new.

\subsection{Provenance and reading boundary}
Repository comparisons use snapshot \snap. We inspected the series group
index, the Hahn--Fuchsian package README, the article's setup and normal-form
section and its further-questions section, the resonance-block package
README, and the well-based-support Lean module \cite{repo-index,repo-hahn,
repo-block,repo-lean}. The complete large canonical volume and every other
research arrival were not audited. This is a targeted comparison, not a
claim of exhaustive novelty checking across the repository or literature.
No repository files were changed.

\section{Fixed supports, directed paths, and coefficient spaces}
\label{sec:setup}

\subsection{The differential system}
Fix $r\geq2$, real numbers $\lambda_1,\ldots,\lambda_r$, and
\[
 \Lambda=\diag(\lambda_1,\ldots,\lambda_r),\qquad
 D=x\frac{d}{dx}.
\]
Let $\G$ be a directed graph on $\{1,\ldots,r\}$ whose edges have the form
$i\to j$ with $i<j$. For every edge $e=(i,j)$ fix a nonempty well-ordered
set $E_e\subset\R_{>0}$. A missing or empty edge is simply deleted.
Consider
\begin{equation}\label{eq:system}
 DY=(\Lambda+A(x))Y,\qquad
 A_{ij}(x)=\sum_{\alpha\in E_{ij}}a_{ij,\alpha}x^\alpha,
\end{equation}
with $A_{ij}=0$ off the prescribed edges. Coefficients are complex and may
vanish. The support sets, graph, and spectrum are fixed; the coefficients
are the variable data.

We work formally in
\[
 \K=\C((x^\R)),\qquad \K[L],\qquad
 D(x^\alpha)=\alpha x^\alpha,\quad DL=1.
\]
Here a Hahn support is well ordered in the ordinary order of real
exponents. The variable $L$ is algebraically independent formally and is
later evaluated as $\log x$. Smaller powers are dominant as $x\to0$.

For a path $p=(i_0,i_1,\ldots,i_m)$ put
\begin{equation}\label{eq:pathsum}
 E_p=E_{i_0i_1}+\cdots+E_{i_{m-1}i_m},\qquad
 \rho_p=\lambda_{i_0}-\lambda_{i_m}.
\end{equation}
The length $m$ is at most $r-1$. Write $S_{ij}$ for the finite union of all
$E_p$ with endpoints $i,j$. No full additive semigroup is needed to describe
these supports.

\begin{lemma}[Finite path support arithmetic]\label{lem:support}
Every $E_p$ and $S_{ij}$ is well ordered. For a fixed path and a fixed real
$s$, there are only finitely many tuples
$(\alpha_1,\ldots,\alpha_m)$, with each $\alpha_k$ in its edge support,
whose sum is $s$.
\end{lemma}
\begin{proof}
For two well-ordered subsets $U,V$ of a linearly ordered abelian group,
$U+V$ is well ordered and every fixed sum has finitely many decompositions.
In the present real case this can also be seen directly. Every sequence in
a well-ordered set has a nondecreasing subsequence. A strictly decreasing
sequence $u_n+v_n$ would therefore, after taking such a subsequence of the
$u_n$, give a strictly decreasing sequence of $v_n$, a contradiction.
If infinitely many distinct pairs had the same sum, their distinct first
coordinates would have a strictly increasing subsequence, forcing strictly
decreasing second coordinates. Induction proves both assertions for any
fixed number of factors. There are only finitely many paths in the graph,
and a finite union of well-ordered subsets of a linear order is well
ordered.
\end{proof}

The elementary subsequence fact used here follows by choosing successive
tail minima; equivalently, a sequence with no nondecreasing subsequence
would yield an infinite strictly decreasing subsequence. The repository
provides the corresponding two-factor support and finite-antidiagonal
bridges as \texttt{Fabius.neumann\_add\_isPWO} and
\texttt{Fabius.neumann\_finite\_decompositions} \cite{repo-lean}.

\subsection{Divided powers and absolute coefficient norms}
Put $\dplog{k}=L^k/k!$. For
$u(L)=\sum_{k=0}^d u_k\dplog{k}$ define
\begin{equation}\label{eq:polynorm}
 |u|_{\mathrm{dp}}=\sum_{k=0}^d|u_k|.
\end{equation}
On a fixed degree bound, differentiation shifts the divided-power
coefficients and has norm at most one. Integration with zero constant
coefficient also has norm at most one, as a map to the next degree.

For $R>0$, a fixed well-ordered $E$, and a positive finite weight
$w:E\to(0,\infty)$, let
\begin{equation}\label{eq:inputnorm}
 X_R(E,w)=\left\{a:\sum_{\alpha\in E}|a_\alpha|R^\alpha
                                      w(\alpha)<\infty\right\}.
\end{equation}
This is a Banach space, isometric to ordinary $\ell^1(E)$ after multiplying
coordinates by $R^\alpha w(\alpha)$. For a fixed output support $S$ and
degree bound $d$, use
\begin{equation}\label{eq:outputnorm}
 \norm{h}_{R,\mathrm{dp}}
   =\sum_{s\in S}R^s|h_s(L)|_{\mathrm{dp}},\qquad
 h=\sum_{s\in S}h_s(L)x^s.
\end{equation}
The resulting space is likewise Banach. Subsets of these fixed supports
remain Hahn supports, so completion does not leave the formal class.

Let $\A_R$ denote the collection of Hahn series with finite norm
$\sum |f_s|R^s$, allowing any individual well-ordered support. It is a
normed algebra under finite coefficient convolution. We do \emph{not}
claim that this unrestricted-support algebra is complete: an $\ell^1$
limit of changing supports need not have well-ordered support. All
closed-graph arguments below take place in the fixed-support Banach spaces
just defined.

Membership in $\A_R[L]$ means that each of finitely many log coefficients
belongs to $\A_R$. Multiplication by a finite sum of real monomials times
log polynomials preserves this membership. Indeed, multiplication by
$x^\beta$ scales the coefficient norm by $R^\beta$, and multiplication of
fixed-degree polynomials changes norms by only a finite combinatorial
constant. No completeness assertion is required for this observation.

\section{The exact path kernel}
\label{sec:kernel}

\subsection{A polynomial inverse including exact resonance}
For $\beta\in\R$ define a linear operator on $\C[L]$ by
\begin{equation}\label{eq:T}
 T_\beta u=
 \begin{cases}
 \displaystyle\sum_{k=0}^{\deg u}
          \frac{(-1)^k}{\beta^{k+1}}\partial_L^k u,&\beta\ne0,\\[6pt]
 \displaystyle\int_0^L u(v)\,dv,&\beta=0.
 \end{cases}
\end{equation}
The zero polynomial is mapped to zero.

\begin{lemma}[Polynomial right inverse]\label{lem:T}
For every $u$, $(\beta+\partial_L)T_\beta u=u$.
If $\beta\ne0$, $T_\beta u$ is the unique polynomial solution. If
$\beta=0$, it is the unique polynomial solution with zero constant
coefficient. For $u\ne0$, a nonzero $\beta$ preserves degree and divides
the divided-power leading coefficient by $\beta$; $T_0$ raises degree by
one and preserves that leading coefficient.
\end{lemma}
\begin{proof}
For nonzero $\beta$ the finite geometric sum telescopes, because a
sufficiently high derivative of $u$ is zero. A nonzero polynomial cannot
lie in the kernel of $\beta+\partial_L$: its highest-degree coefficient
would be multiplied by the nonzero scalar $\beta$. For $\beta=0$,
differentiation and zero-constant integration are inverse in the stated
sense. The leading-coefficient assertions follow directly in the
$\dplog{k}$ basis.
\end{proof}

For degree at most $d$ and $\beta\ne0$,
\begin{equation}\label{eq:Tbound}
 |T_\beta u|_{\mathrm{dp}}
 \leq\left(\sum_{k=0}^d|\beta|^{-k-1}\right)|u|_{\mathrm{dp}}.
\end{equation}
In particular, a uniform lower bound on nonzero $|\beta|$ bounds all
operators needed for paths of bounded length.

\subsection{Suffix spectral differences}
Fix a path $p=(i_0,\ldots,i_m)$ and an exponent tuple
$\boldsymbol\alpha=(\alpha_1,\ldots,\alpha_m)$. Define
\begin{equation}\label{eq:betas}
 \beta_k=\alpha_k+\cdots+\alpha_m
               -(\lambda_{i_{k-1}}-\lambda_{i_m}),
 \qquad 1\leq k\leq m,
\end{equation}
and the polynomial path kernel
\begin{equation}\label{eq:kernel}
 Q_p(\boldsymbol\alpha;L)
       =T_{\beta_1}T_{\beta_2}\cdots T_{\beta_m}1.
\end{equation}
The rightmost operator is applied first. The spectral differences refer
to the \emph{same terminal vertex}. Replacing them by the individual edge
differences is generally incorrect.

\begin{lemma}[Degree and nonvanishing of a path kernel]\label{lem:leading}
Let $z$ be the number of zero entries among $\beta_1,\ldots,\beta_m$.
Then
\begin{equation}\label{eq:leading}
 \deg_L Q_p=z,\qquad
 \lc Q_p=\prod_{\beta_k\ne0}\beta_k^{-1}.
\end{equation}
In particular, no path kernel is the zero polynomial.
\end{lemma}
\begin{proof}
Start with the polynomial $1$ and apply Lemma~\ref{lem:T} from right to
left. Every exact resonance raises degree by one without changing the
divided-power leading coefficient. Every other step multiplies that
coefficient by $\beta_k^{-1}$. These factors are nonzero.
\end{proof}

This formula is useful for the necessity direction. It avoids attempting
to estimate a potentially cancellation-prone constant coefficient of a
log polynomial. The highest divided-power coefficient is a single product.

\section{A canonical Frobenius fundamental matrix}
\label{sec:formal}

We first construct a polynomial-logarithmic factor $H$ such that
$Y=Hx^\Lambda$ solves \eqref{eq:system}. The log-free normalizer is denoted
by $P$ and constructed separately in Section~\ref{sec:gauge}; keeping the
two letters distinct prevents a change of normalization from being hidden.

\begin{theorem}[Canonical path expansion]\label{thm:formal}
There is a unique upper-unitriangular matrix $H\in\Mat_r(\K[L])$ with
positive off-diagonal support that satisfies
\begin{equation}\label{eq:H-equation}
 DH=[\Lambda,H]+AH
\end{equation}
and the following normalization: whenever
$s=\lambda_i-\lambda_j>0$, the polynomial coefficient $[x^s]H_{ij}$ has
zero constant coefficient. Its entries are
\begin{equation}\label{eq:H-path}
 H_{ij}=
 \sum_{p\in\Paths(i,j)}\ 
 \sum_{\alpha_k\in E_{i_{k-1}i_k}}
 \left(\prod_{k=1}^m a_{i_{k-1}i_k,\alpha_k}\right)
 x^{\alpha_1+\cdots+\alpha_m}Q_p(\boldsymbol\alpha;L)
 \quad(i<j).
\end{equation}
Every coefficient sum is finite, $\supp H_{ij}\subseteq S_{ij}$, and
$\deg_LH_{ij}\leq j-i\leq r-1$. Thus $Hx^\Lambda$ is a formal
fundamental matrix.
\end{theorem}
\begin{proof}
Proceed by increasing $j-i$. With $H_{jj}=1$, the $(i,j)$ equation is
\begin{equation}\label{eq:H-recursion}
 (D-\rho_{ij})H_{ij}
       =A_{ij}+\sum_{i<k<j}A_{ik}H_{kj},\qquad
 \rho_{ij}=\lambda_i-\lambda_j.
\end{equation}
Its right side has already been defined. On a coefficient $x^s u(L)$,
the left side acts as $x^s(s-\rho_{ij}+\partial_L)u$. Apply
$T_{s-\rho_{ij}}$ to each coefficient. Lemma~\ref{lem:T} gives exactly the
stated normalization and uniqueness. Positivity of input exponents and
Lemma~\ref{lem:support} justify the coefficient operations.

Expanding the recursion prepends one edge to a suffix path. The
homological factor at that step is the corresponding $\beta_k$ in
\eqref{eq:betas}. Induction therefore gives \eqref{eq:H-path}. There are
at most $j-i$ edges in a path, and each $T_0$ raises degree by at most one.
The determinant of $H$ is one, so $Hx^\Lambda$ is invertible.
\end{proof}

The theorem requires neither analytic convergence nor local finiteness of
the supports. Infinite sets below a finite spectral value do not invalidate
the formal recursion: a fixed coefficient still has only finitely many
contributing tuples.

\section{An exact weighted convergence criterion}
\label{sec:weighted}

\subsection{The path norm}
Assign a weight $w_e:E_e\to(0,\infty)$ to each edge. For a path $p$ define
\begin{equation}\label{eq:Mp}
 M_p(w)=\sup_{\boldsymbol\alpha\in\prod E_e}
 \frac{|Q_p(\boldsymbol\alpha;L)|_{\mathrm{dp}}}
      {\prod_{e\in p}w_e(\alpha_e)}\ \in[0,\infty].
\end{equation}
The exponent-radius factors do not occur: $R^{\sum\alpha_e}$ cancels
against $\prod R^{\alpha_e}$.

\begin{theorem}[Exact independent-input criterion]\label{thm:weighted}
Fix the graph, supports, spectrum, weights, and $R>0$. The following are
equivalent:
\begin{enumerate}[label=(\roman*)]
\item Every independent choice $a_e\in X_R(E_e,w_e)$ gives a canonical
      $H$ with all coefficients in $\A_R[L]$.
\item $M_p(w)<\infty$ for every directed path $p$.
\end{enumerate}
If these conditions hold, then
\begin{equation}\label{eq:weightedbound}
 \norm{H_{ij}}_{R,\mathrm{dp}}
 \leq\sum_{p\in\Paths(i,j)}M_p(w)
                \prod_{e\in p}\norm{a_e}_{X_R(E_e,w_e)}.
\end{equation}
Moreover, $M_p(w)$ is exactly the multilinear operator norm of the isolated
path contribution, not merely an upper bound.
\end{theorem}

\begin{proof}[Sufficiency and exact norm]
For a fixed path, the triangle inequality, \eqref{eq:Mp}, and the identity
$R^{\sum\alpha_e}=\prod R^{\alpha_e}$ bound the output norm by
$M_p(w)$ times the product of the input norms. Absolute summability
justifies the resulting product of scalar sums. Adding the finitely many
paths proves \eqref{eq:weightedbound}.

To see sharpness, give each edge a single nonzero coefficient at
$\alpha_e$, with magnitude $(R^{\alpha_e}w_e(\alpha_e))^{-1}$.
Every input has norm one and the output norm is exactly the ratio in
\eqref{eq:Mp}. Taking the supremum proves the reverse norm inequality.
\end{proof}

The necessity direction needs more than observing that single-coefficient
examples have arbitrarily large norms. A map can be unbounded on a sequence
of finite inputs without that observation alone producing one infinite
input with divergent output. The following standard functional-analytic
argument supplies the missing quantifier.

\begin{lemma}[Automatic boundedness for the path map]\label{lem:closedgraph}
Suppose the formal multilinear contribution of a fixed path maps every
tuple of its fixed-support Banach input spaces into the fixed-support
Banach output space of degree at most $r-1$. Then that multilinear map is
bounded.
\end{lemma}
\begin{proof}
Fix all but one input. Consider a sequence in the remaining input space
converging in norm, whose outputs also converge in output norm. Every
coefficient of a fixed exponent and fixed divided power is a finite sum
of input coordinates by Lemma~\ref{lem:support}. Coordinate evaluation is
continuous in each weighted $\ell^1$ space. Thus the output limit has the
same coordinates as the formal output of the input limit. The graph of
the resulting everywhere-defined linear map is closed. The closed graph
theorem makes it continuous.

The multilinear map is therefore separately continuous. Separate
continuity on a finite product of Banach spaces implies joint boundedness.
For completeness, argue by induction on the number of factors. Fixing the
first factor, the induction hypothesis bounds the map on the unit balls
of the remaining factors. Regard those remaining arguments as indexing a
family of bounded linear maps of the first factor. That family is
pointwise bounded by the induction hypothesis. The uniform boundedness
principle supplies a uniform operator bound, as required.
\end{proof}

\begin{proof}[Necessity in Theorem~\ref{thm:weighted}]
Fix a path and set every edge not on that path equal to zero. Between its
endpoints, the retained graph has just that full path. Assumption (i)
therefore says that its formal multilinear contribution maps every tuple
of edge inputs into the fixed-support output Banach space. By
Lemma~\ref{lem:closedgraph} it has finite norm. The single-coefficient tests
in the first part of the proof show that this norm is $M_p(w)$.
\end{proof}

\begin{remark}[What the word ``exact'' means]
The criterion is exact for the specified independent product of weighted
$\ell^1$ spaces and for absolute coefficient convergence. It is not a
necessary test for one selected input with cancellations, nor for a family
whose edge coefficients satisfy additional algebraic relations. Nor does it
claim that the infinite supremum is computable from every possible
presentation of the supports and weights.
\end{remark}

\section{The path--spectrum gap theorem}
\label{sec:gap}

For every path define its nonresonant gap
\begin{equation}\label{eq:gap}
 \eta_p=\inf\{|s-\rho_p|:s\in E_p,\ s\ne\rho_p\},
\end{equation}
with infimum $+\infty$ for an empty set. Exact resonance is explicitly
removed from the infimum.

\begin{theorem}[Sharp accessible-spectrum criterion]\label{thm:gap}
For the system \eqref{eq:system}, the following conditions are equivalent:
\begin{enumerate}[label=(\roman*)]
\item For one fixed $R>0$, every input with
      $\sum_{e,\alpha}|a_{e,\alpha}|R^\alpha<\infty$ has
      $H\in\Mat_r(\A_R[L])$.
\item The same assertion holds for every $R>0$.
\item $\eta_p>0$ for every directed path $p$.
\item No $\rho_p$ is a left accumulation point of its own path sumset
      $E_p$.
\end{enumerate}
Thus a spectral difference matters only for paths with its specified
endpoints. It is not enough to take an undifferentiated union of all path
sumsets and all spectral differences.
\end{theorem}

\begin{proof}
Assume (iii). There are finitely many paths, so every nonzero $\beta_k$
in every path tuple is bounded away from zero by one common $\eta>0$.
Indeed, a suffix is itself a path, and its $\beta_k$ is exactly a
nonresonant endpoint difference for that suffix. Let $d=r-1$ and
$h=\min(1,\eta)$. Equations~\eqref{eq:Tbound} and
$|T_0u|_{\mathrm{dp}}=|u|_{\mathrm{dp}}$ give the uniform bound
\[
 |Q_p|_{\mathrm{dp}}
 \leq\bigl((d+1)h^{-(d+1)}\bigr)^{|p|}.
\]
Consequently every unweighted $M_p$ is finite. Theorem~\ref{thm:weighted}
proves (ii), which immediately implies (i).

Conversely, suppose $\eta_p=0$ for a path of length $m$. Choose tuples
with sums $s_n\ne\rho_p$ and $s_n\to\rho_p$. Their first homological
factors satisfy $\beta_{1,n}=s_n-\rho_p\to0$. Because all input exponents
are positive and $s_n$ is bounded, every suffix sum is bounded. Hence there
is a fixed $B\geq1$ such that every $|\beta_{k,n}|\leq B$. By
Lemma~\ref{lem:leading}, the highest divided-power coefficient of the
kernel satisfies
\begin{equation}\label{eq:blowup}
 |Q_p(\boldsymbol\alpha_n)|_{\mathrm{dp}}
 \geq\frac{1}{B^{m-1}|s_n-\rho_p|}\longrightarrow\infty.
\end{equation}
Zero suffix factors cause no problem: they are omitted from the leading
coefficient product. Thus $M_p=\infty$, contradicting (i) by
Theorem~\ref{thm:weighted}. This proves (i)$\Rightarrow$(iii).

Finally, a well-ordered real set cannot accumulate at a finite point from
above: one could successively choose smaller elements approaching that
point and obtain an infinite decreasing sequence. If its distance from a
point, with exact equality removed, is zero, it therefore accumulates from
below. The reverse implication is immediate. Apply this to each $E_p$ to
obtain (iii)$\Leftrightarrow$(iv).
\end{proof}

\begin{corollary}[Small inputs; no rescue by reducing the radius]
\label{cor:allradii}
If the gap condition fails, then for every $\eps>0$ there is an input
supported on a single directed chain, with all input exponents in a
bounded interval and
$\sum_{e,\alpha}|a_{e,\alpha}|<\eps$, whose canonical $H$ fails absolute
coefficient convergence at every radius $R>0$.
\end{corollary}
\begin{proof}
Use a bad path from the proof of Theorem~\ref{thm:gap}. All exponents in
its witnessing tuples lie below one finite $B_0$. Restrict each edge to
$E_e\cap(0,B_0]$. The corresponding path norm remains infinite, so the
necessity argument in Theorem~\ref{thm:weighted}, now at radius one,
provides a tuple of $\ell^1$ inputs whose path output is not in $\ell^1$.
Set all other edges to zero. Every output exponent lies in
$(0,(r-1)B_0]$. On a bounded interval, the weights $R^s$ are bounded above
and below by positive constants, so nonconvergence persists at every
radius. Scaling all retained inputs by one sufficiently small positive
number scales this isolated path contribution by a nonzero power of that
number and makes the total input norm less than $\eps$.
\end{proof}

For $H$, the conclusion refers to its polynomial-logarithmic coefficient
norm, not to every possible exceptional numerical value of $L$. The next
section transfers it to a genuinely log-free gauge, where absolute
convergence at radius $R$ is the ordinary absolute generalized-series test.

\section{Equivalence with the normalized log-free gauge}
\label{sec:gauge}

\subsection{Triangular normal form without an analytic assumption}
For $i<j$ put $\rho_{ij}=\lambda_i-\lambda_j$. We seek an upper-unitriangular
log-free matrix $P$ and a strictly upper-triangular resonant matrix
$\mathcal R$ satisfying
\begin{equation}\label{eq:P-equation}
 DP=(\Lambda+A)P-P(\Lambda+\mathcal R),
\end{equation}
where
\begin{equation}\label{eq:Rform}
 \mathcal R_{ij}=b_{ij}x^{\rho_{ij}}\quad\text{if }\rho_{ij}>0,
 \qquad \mathcal R_{ij}=0\quad\text{otherwise},
\end{equation}
and $[x^{\rho_{ij}}]P_{ij}=0$ whenever $\rho_{ij}>0$.

\begin{proposition}[Unique normalized log-free pair]\label{prop:normal}
There is a unique pair $(P,\mathcal R)$ of this form with positive
 off-diagonal support. The support of each entry is contained in its path
sumset union. The matrix $B=(b_{ij})$ is constant and nilpotent, and
\begin{equation}\label{eq:normalfund}
 F=P x^\Lambda e^{BL}
\end{equation}
is a formal fundamental matrix. Every coefficient of $B$ depends
polynomially on finitely many original input coefficients, with constants
determined by the fixed support and spectral data.
\end{proposition}
\begin{proof}
Proceed again by increasing $j-i$. Define, using previously computed
entries,
\begin{equation}\label{eq:Fij}
 G_{ij}=A_{ij}+\sum_{i<k<j}A_{ik}P_{kj}
                       -\sum_{i<k<j}P_{ik}\mathcal R_{kj}.
\end{equation}
If $\rho_{ij}>0$, put
$b_{ij}=[x^{\rho_{ij}}]G_{ij}$; otherwise put $b_{ij}=0$. Set
\begin{equation}\label{eq:Pcoeff}
 P_{ij}=\sum_{s\ne\rho_{ij}}
                 \frac{[x^s]G_{ij}}{s-\rho_{ij}}x^s.
\end{equation}
All supports are positive, so there is no omitted nonpositive resonance.
This is exactly the coefficient equation for
$(D-\rho_{ij})P_{ij}=G_{ij}-\mathcal R_{ij}$ and proves existence and
uniqueness.

Every product in \eqref{eq:Fij} concatenates two permitted paths.
Consequently its support remains within $S_{ij}$. Lemma~\ref{lem:support}
then proves finite coefficient dependence at each fixed resonance.
Recursion uses only additions, multiplications, and division by fixed
nonzero real numbers, so the dependence on input coefficients is
polynomial. There are finitely many matrix entries.

Since $B$ is strictly upper triangular, $B^r=0$. Equation~\eqref{eq:Rform}
is $\mathcal R=x^\Lambda Bx^{-\Lambda}$. Differentiating
\eqref{eq:normalfund} and using \eqref{eq:P-equation} now proves the
fundamental-matrix assertion.
\end{proof}

The finite-dependence statement is compatible with, and not a replacement
for, the more general finite-ancestry result in the predecessor
\cite{repo-hahn}. Our triangular construction is included to make the
normalization comparison self-contained.

\subsection{Why the two convergence questions are the same}
\begin{lemma}[Constants]\label{lem:constants}
The constants of $\K[L]$ under $D$ are exactly $\C$.
\end{lemma}
\begin{proof}
A constant Hahn series has only exponent zero, since its coefficient at
$x^s$ is multiplied by $s$ under $D$. Also, no derivative of a Hahn series
has a nonzero $x^0$ coefficient. If a nonzero polynomial
$\sum_{k=0}^d f_kL^k$ had derivative zero and $d>0$, its highest coefficient
would satisfy $Df_d=0$ and hence be a nonzero complex constant. The next
equation would be $Df_{d-1}=-df_d$, contradicting the missing $x^0$
coefficient. Thus $d=0$.
\end{proof}

\begin{theorem}[Convergence is independent of these two normalizations]
\label{thm:gauge-equivalence}
For every fixed formal input and every $R>0$,
\begin{equation}\label{eq:equivalence}
 H\in\Mat_r(\A_R[L])
 \quad\Longleftrightarrow\quad P\in\Mat_r(\A_R).
\end{equation}
Consequently Theorem~\ref{thm:gap} is an exact universal convergence
criterion for the normalized log-free gauge $P$. The weighted membership
criterion of Theorem~\ref{thm:weighted} also transfers to $P$ for the same
weighted input classes. If the gap condition fails, the small input in
Corollary~\ref{cor:allradii} has a log-free normalizer diverging absolutely
at every positive radius.
\end{theorem}
\begin{proof}
Both $F=P x^\Lambda e^{BL}$ and $Y=Hx^\Lambda$ are invertible fundamental
matrices over $\K[L]$. Their inverses lie in the same ring: upper
unitriangular inverses and $e^{-BL}$ are finite sums, and the diagonal
monomials are invertible. Therefore
$C=F^{-1}Y$ has derivative zero and belongs to $\mathrm{GL}_r(\C)$ by
Lemma~\ref{lem:constants}. We obtain
\begin{equation}\label{eq:factorcomparison}
 H=P\,x^\Lambda e^{BL}C x^{-\Lambda},\qquad
 P=H\,x^\Lambda C^{-1}e^{-BL}x^{-\Lambda}.
\end{equation}
Each rightmost factor is a finite sum of monomials times log polynomials.
Multiplication by it preserves $\A_R[L]$ membership. Since $P$ is
log-free, \eqref{eq:equivalence} follows. The universal and all-radii
conclusions are then immediate.
\end{proof}

This is an equivalence of membership for each input. It does not identify
the path norm of $H$ with an identical operator norm of $P$, because the
finite factors $B$ and $C$ depend on the input. The exact norm assertion in
Theorem~\ref{thm:weighted} remains an assertion about the isolated
polynomial-logarithmic path operator.

On every smaller radius $0<R'<R$, an absolutely convergent gauge evaluates
to an analytic function on a chosen logarithmic cover. Differentiation is
legitimate there: multiplying coefficients by $s^k$ is controlled by the
exponential margin $(R'/R)^s$, while all supports under consideration are
bounded below. The formal equations therefore become analytic equations.
No differentiability assertion at the outer boundary $|x|=R$ is needed.

\section{What three and four dimensions already show}
\label{sec:examples}

\subsection{A bad semigroup accumulation that cannot occur in a matrix word}
\begin{example}[The whole-semigroup condition is not necessary]
\label{ex:harmless}
Take
\[
 \Lambda=\diag(3,1,0),\qquad
 E_{12}=\{3/2-1/n:n\geq2\},\qquad E_{23}=\{10\},
\]
with no other edges. The three permitted paths have data
\[
 \begin{array}{c|c|c}
 p&E_p&\rho_p\\ \hline
 1\to2&E_{12}&2\\
 2\to3&\{10\}&1\\
 1\to2\to3&E_{12}+10&3.
 \end{array}
\]
Their gaps are at least $1/2$, $9$, and $8$, respectively. Thus every
absolutely convergent input has an absolutely convergent normalized gauge.
Yet the full semigroup generated by $E_{12}\cup\{10\}$ contains
$3-2/n\to3$ from below. The positive spectral difference $3$ is therefore
a left accumulation point of that semigroup.
\end{example}

Here $2E_{12}$ describes a forbidden word, because $E_{12}E_{12}=0$ as a
matrix product. Writing $A_{12}=\sum a_\alpha x^\alpha$ and
$A_{23}=b x^{10}$, the actual nonresonant entries are
\[
 P_{12}=\sum\frac{a_\alpha}{\alpha-2}x^\alpha,\qquad
 P_{23}=\frac b9x^{10},\qquad
 P_{13}=\sum\frac{a_\alpha b}{9(\alpha+7)}x^{\alpha+10}.
\]
Every denominator is uniformly safe. This answers the necessity question
negatively if it is phrased using the full generated semigroup without
matrix-path information.

\subsection{Safe edges whose composition is unsafe}
\begin{example}[An accumulating two-letter sumset]\label{ex:hidden}
Take
\[
 \Lambda=\diag(3,2,0),\qquad A_{12}=x^{1/2},\qquad
 A_{23}=\sum_{n\geq3}n^{-p}x^{5/2-1/n},\qquad p>1.
\]
The individual gaps are $1/2$ and at least $1/6$. Nevertheless,
$1/2+(5/2-1/n)=3-1/n\to3$, the endpoint difference for the composed path.
There are no exact resonances, and
\begin{align}
 P_{12}&=-2x^{1/2},\nonumber\\
 P_{23}&=\sum_{n\geq3}\frac{n^{-p}}{1/2-1/n}x^{5/2-1/n},\label{eq:hidden}\\
 P_{13}&=-\sum_{n\geq3}\frac{n^{1-p}}{1/2-1/n}x^{3-1/n}.\nonumber
\end{align}
Thus the original input and both edge normalizations converge for $p>1$,
whereas the full normalizer converges exactly when $p>2$.
\end{example}

The equivalence at every radius follows because all displayed exponent
sets are bounded. The coefficients of $P_{13}$ have one sign, so there is
no cancellation issue. This is the specific two-letter phenomenon proposed
in the predecessor's further-questions section, now placed inside a
necessary-and-sufficient theorem for all acyclic patterns.

\subsection{Why independent inputs cannot be silently replaced by coupled inputs}
\label{sec:cancellation}
Take the same $f(x)=\sum_{n\geq3}n^{-p}x^{5/2-1/n}$, and set
\[
 \Lambda=\diag(3,2,2,0),\qquad
 A_{12}=A_{13}=x^{1/2},\qquad A_{24}=f,\quad A_{34}=-f,
\]
with all other entries zero. The two paths from $1$ to $4$ have identical
kernels and opposite coefficient products. Hence $H_{14}=P_{14}=0$.
The remaining nonzero entries have the safe one-edge denominators in
Example~\ref{ex:hidden}, and the whole system normalizes convergently for
$p>1$ even though each of the two path sumsets accumulates at $3$.

This does not contradict Theorem~\ref{thm:gap}: the selected coefficients
are correlated. In matrix-direction language,
$(E_{12}+E_{13})(E_{24}-E_{34})=0$. A graph that records the four edges
but forgets this relation loses essential information. The exact weighted
criterion handles all independent inputs; a genuinely optimal criterion
for constrained matrix directions must retain the corresponding algebraic
relations among words.

\section{Resonant chains and sharp arbitrary-order moment loss}
\label{sec:chain}

\subsection{The family}
For $r\geq2$, let
\begin{equation}\label{eq:chain}
 \lambda_i=r-i,\qquad
 A_{12}=\sum_{n\geq2}c_nx^{1-1/n},\qquad
 A_{j,j+1}=x\quad(2\leq j<r),
\end{equation}
with no other edges. The first edge approaches its resonance from below;
every downstream edge is exactly resonant. Put $q=r-2$.

\begin{theorem}[Exact moment threshold]\label{thm:chain}
For every $R>0$, the input in \eqref{eq:chain} is absolutely convergent
exactly when $\sum|c_n|<\infty$. Its normalized log-free gauge is
\begin{equation}\label{eq:chainP}
 P_{1j}=-\sum_{n\geq2}c_n n^{j-1}x^{j-1-1/n}
       \quad(2\leq j\leq r),
\end{equation}
with all other off-diagonal entries zero. Its resonant constant matrix is
$B_{j,j+1}=1$ for $2\leq j<r$, with all other entries zero. The full gauge
converges absolutely if and only if
\begin{equation}\label{eq:moment}
 \sum_{n\geq2}|c_n|n^{r-1}<\infty.
\end{equation}
In particular, for $c_n=n^{-p}$ the input converges for $p>1$, while the
normalizer converges exactly for $p>r$.
\end{theorem}
\begin{proof}
For $P_{12}$, division by $(1-1/n)-1=-1/n$ gives
\eqref{eq:chainP}. The downstream input edges are placed entirely in
$\mathcal R$, and their gauge entries are zero. For $j>2$, the only
nonzero term in \eqref{eq:Fij} is
$-P_{1,j-1}\mathcal R_{j-1,j}=-xP_{1,j-1}$.
Dividing again by $-1/n$ multiplies the previous coefficient by $n$.
Induction proves the formula and the stated $B$.

For fixed $j$ and $R>0$, the factors $R^{j-1-1/n}$ are bounded above and
below by positive constants. Thus convergence of $P_{1j}$ is equivalent
to $\sum|c_n|n^{j-1}<\infty$. The strongest condition occurs at $j=r$.
The same bounded-exponent observation proves the input assertion. Finally,
$\sum n^{r-1-p}$ converges exactly when $p>r$.
\end{proof}

The loss can therefore be made arbitrarily large by increasing the
nilpotent chain length. A single division by the first small denominator
does not measure the full cost of normalization.

\subsection{The logarithmic explanation}
For $2\leq j\leq r$,
\[
 H_{jr}=x^{r-j}\frac{L^{r-j}}{(r-j)!}.
\]
For a fixed $n$ in the first edge, the remaining $q$ homological factors
are zero, so the top-right kernel is
\begin{equation}\label{eq:chainQ}
 T_{-1/n}\dplog{q}
       =-\sum_{k=0}^q n^{k+1}\dplog{q-k}.
\end{equation}
Consequently
\begin{equation}\label{eq:chainH}
 H_{1r}=-\sum_{n\geq2}c_nx^{q+1-1/n}
                      \sum_{k=0}^q n^{k+1}\frac{L^{q-k}}{(q-k)!}.
\end{equation}
Its divided-power coefficient norm is exactly
$\sum |c_n|R^{q+1-1/n}\sum_{k=0}^q n^{k+1}$, confirming the same moment
condition.

For independent weighted inputs on the first edge, with the downstream
coefficients allowed to vary over their one-dimensional spaces, the
minimal weight growth is therefore $w(1-1/n)\asymp n^{r-1}$: more precisely,
universal convergence holds exactly when
$\sup_{n\geq2} n^{r-1}/w(1-1/n)<\infty$. This is an exact consequence of
the path norm, not a guessed submultiplicative majorant.

\section{Stable analytic realization beyond the moment threshold}
\label{sec:analytic}

\subsection{A convergent kernel instead of divergent constants}
Assume only $\sum_{n\geq2}|c_n|<\infty$. Put $t=\log(1/x)$, initially
for $0<x\leq1$. Define the entire function
\begin{equation}\label{eq:Psi}
 \Psi_q(z)=\frac1{q!}\int_0^1 e^{zu}u^q\,du
          =\frac1{q!}\sum_{k\geq0}\frac{z^k}{k!(q+k+1)}.
\end{equation}
The integral or the positive series is a stable way to evaluate it near
$z=0$.

\begin{theorem}[Basepoint-normalized realization]\label{thm:analytic-chain}
The chain \eqref{eq:chain} has a unique analytic fundamental matrix on the
logarithmic cover with $Y(1)=I$. Writing it as
$Y^{\mathrm{bp}}=H^{\mathrm{bp}}x^\Lambda$, its top-right entry is
\begin{align}
 H^{\mathrm{bp}}_{1r}(e^{-t})
 &=\frac{(-1)^{q+1}e^{-(q+1)t}}{q!}
      \sum_{n\geq2}c_n\int_0^t e^{v/n}v^q\,dv\label{eq:bp-int}\\
 &=(-1)^{q+1}e^{-(q+1)t}t^{q+1}
      \sum_{n\geq2}c_n\Psi_q(t/n).\label{eq:bp-stable}
\end{align}
The analogous formula with $q=j-2$ gives $H^{\mathrm{bp}}_{1j}$.
All downstream entries are the finite resonant expressions already given.
These functions are entire in $t$ and satisfy the differential system.

Moreover, this analytic matrix realizes every fixed finite exponent-block
prefix of the canonical expansion, even when \eqref{eq:moment} fails.
\end{theorem}

\begin{proof}
For downstream entries, direct integration gives
$x^{k-j}L^{k-j}/(k-j)!$. For the top-right entry, variation of constants
in \eqref{eq:H-recursion}, with value zero at $x=1$, gives
\[
 H^{\mathrm{bp}}_{1r}(x)
 =x^{q+1}\sum_{n\geq2}c_n\int_0^{\log x}
                      e^{-u/n}\frac{u^q}{q!}\,du.
\]
The substitution $u=-v$ proves \eqref{eq:bp-int}, and $v=tu$ gives
\eqref{eq:bp-stable}. For complex $t$ in a compact set,
$|\Psi_q(t/n)|\leq e^{|t|/2}/(q+1)!$. Hence the sum converges normally
under $\sum|c_n|<\infty$, and the same holds after any fixed number of
$t$ derivatives. The equations and the basepoint normalization follow
termwise. Triangular variation of constants also proves uniqueness.

It remains to justify the formal-prefix assertion. For each fixed $n$,
elementary integration gives, with $E_q(z)=\sum_{j=0}^qz^j/j!$,
\begin{equation}\label{eq:subtraction}
 H^{\mathrm{bp}}_{1r,n}(e^{-t})
 =c_ne^{-(q+1)t}n^{q+1}
       \bigl[1-e^{t/n}E_q(-t/n)\bigr].
\end{equation}
The term containing $e^{t/n}E_q(-t/n)$ is exactly the $n$th canonical
block in \eqref{eq:chainH}; the remaining term has exponent $q+1$.
Fix $M\geq2$ and subtract the canonical blocks $2\leq n\leq M$.
The finitely many added exponent-$(q+1)$ terms are $O(x^{q+1})$.
For the remaining basepoint terms and $t\geq1$,
\begin{align*}
 \sum_{n>M}|H^{\mathrm{bp}}_{1r,n}(e^{-t})|
 &\leq \frac{e^{-(q+1)t}}{q!}
     \sum_{n>M}|c_n|\int_0^t e^{v/(M+1)}v^q\,dv\\
 &\leq \frac{M+1}{q!}\norm{c}_{\ell^1}
          e^{-(q+1-1/(M+1))t}t^q.
\end{align*}
Thus the remainder is
$O(x^{q+1-1/(M+1)}t^q)$. Its exponent is strictly larger than the last
retained exponent $q+1-1/M$, so it is smaller than that block's monomial
times every fixed inverse power of $t$. The same proof applies to each
top-row entry. Downstream expansions are finite.
\end{proof}

Formula~\eqref{eq:subtraction} must be kept grouped when the relevant
moments diverge. Splitting it into two separately summed expressions would
introduce the divergent quantity $\sum c_nn^{q+1}$. The stable expression
\eqref{eq:bp-stable} never forms that quantity. Indeed,
\[
 1-e^zE_q(-z)=(-1)^{q+1}z^{q+1}\Psi_q(z),
\]
so precisely $q+1$ powers of the small parameter cancel.

This is a selected analytic realization determined by a basepoint. It is
not an assertion that the original canonical Hahn sum converges, or that
formal series without a chosen analytic normalization determine unique
functions. When the strong moment condition holds, the canonical and
basepoint fundamental matrices are related by an ordinary finite constant
matrix. Below that condition, their formal comparison cannot be justified
by summing the divergent constants separately.

\section{A uniform accumulation crossover with an error certificate}
\label{sec:crossover}

Set $c_n=n^{-p}$ with $p>1$. For $N\geq2$ define the grouped basepoint tail
\begin{equation}\label{eq:tail}
 \mathcal U_N(t)=(-1)^{q+1}e^{-(q+1)t}t^{q+1}
                    \sum_{n\geq N}n^{-p}\Psi_q(t/n).
\end{equation}
This is a tail of the convergent grouped representation, not a tail
obtained by separately summing the divergent canonical coefficients.
Let
\begin{equation}\label{eq:gI}
 g_{p,q}(v)=v^{-p}\Psi_q(1/v),\qquad
 I_{p,q}(c)=\int_c^\infty g_{p,q}(v)\,dv.
\end{equation}
Both are positive for $v,c>0$.

\begin{theorem}[Uniform crossover and explicit two-sided error]
\label{thm:crossover}
Let $t>0$, $N\geq2$, $c=N/t$, and
\[
 \mathcal E_{N,t}=(-1)^{q+1}e^{(q+1)t}t^{p-q-2}\mathcal U_N(t).
\]
Then
\begin{equation}\label{eq:rectangle}
 I_{p,q}(c)\leq\mathcal E_{N,t}
       \leq I_{p,q}(c)+\frac{g_{p,q}(c)}{t}.
\end{equation}
The sharper centered estimate is
\begin{equation}\label{eq:trapezoid}
 0\leq\mathcal E_{N,t}-I_{p,q}(c)-\frac{g_{p,q}(c)}{2t}
       \leq\frac{-g'_{p,q}(c)}{8t^2}.
\end{equation}
Consequently, if $N/t\to c_0\in(0,\infty)$,
\[
 \mathcal U_N(t)
 \sim(-1)^{q+1}e^{-(q+1)t}t^{q+2-p}I_{p,q}(c_0).
\]
The error estimates are uniform for $c$ in compact subsets of $(0,\infty)$
and $p$ in compact subsets of $(1,\infty)$, with fixed $q$.
\end{theorem}

\begin{proof}
A direct rescaling gives
\begin{equation}\label{eq:riemann}
 \mathcal E_{N,t}=\frac1t\sum_{n\geq N}g_{p,q}(n/t).
\end{equation}
Expanding \eqref{eq:Psi},
\begin{equation}\label{eq:gseries}
 g_{p,q}(v)=\frac1{q!}\sum_{k\geq0}
                       \frac{v^{-p-k}}{k!(q+k+1)}.
\end{equation}
It is positive, decreasing, convex, and tends to zero. The same series
shows that its derivative tends to zero and that its second derivative is
integrable on every $[c,\infty)$. The ordinary decreasing-function
rectangle bounds for \eqref{eq:riemann} give \eqref{eq:rectangle}.

For the sharper estimate, set $h=1/t$. The trapezoidal error on
$[a,a+h]$ is
\[
 \frac h2\bigl(g(a)+g(a+h)\bigr)-\int_a^{a+h}g(v)\,dv
 =\frac12\int_a^{a+h}(v-a)(a+h-v)g''(v)\,dv.
\]
Convexity makes it nonnegative, and the factor in front of $g''$ is at
most $h^2/8$. Sum over all intervals beginning at $c$.
Their total error is at most
$\frac{h^2}{8}\int_c^\infty g''(v)\,dv=-h^2g'(c)/8$.
The left-endpoint sum in \eqref{eq:riemann} is the infinite trapezoidal sum
plus $hg(c)/2$. This proves \eqref{eq:trapezoid}. Uniformity follows from
the displayed positive series and their derivatives on the stated compact
parameter sets.
\end{proof}

\subsection{Evaluation formulas and the critical scaling}
Termwise integration in \eqref{eq:gseries} gives a useful convergent
representation
\begin{equation}\label{eq:Iseries}
 I_{p,q}(c)=\frac1{q!}\sum_{k\geq0}
       \frac{c^{1-p-k}}{k!(q+k+1)(p+k-1)}.
\end{equation}
For the discrete quantity,
\begin{equation}\label{eq:zeta}
 \mathcal E_{N,t}
 =\frac{t^{p-1}}{q!}\sum_{k\geq0}
       \frac{t^k\zeta(p+k,N)}{k!(q+k+1)},
\end{equation}
where $\zeta(s,N)=\sum_{n\geq N}n^{-s}$ is the Hurwitz zeta tail for
real $s>1$. All terms are positive, and successive-term ratio bounds
$t/(N(k+1))$ and $1/(c(k+1))$ give elementary geometric tail estimates for
\eqref{eq:zeta} and \eqref{eq:Iseries} once those ratios are below one.

The scale $N\asymp t$ has a precise interpretation. For a fixed exponent
index, $t/n$ becomes large, and the canonical exponential block resolves
that term. For indices proportional to $t$, infinitely many nearly
resonant blocks contribute together at scale
$e^{-(q+1)t}t^{q+2-p}$. The moment threshold is $p=q+2=r$; at this threshold
the extra power of $t$ is zero. This does not remove the divergence of the
canonical gauge, which at equality has a harmonic coefficient sum.
The grouped tail, rather than the raw canonical sum, has the crossover
limit.

\subsection{Numerical illustrations}
Table~\ref{tab:crossover} uses $p=5/2$ and $c=1$. It includes dimensions in
which the canonical gauge is divergent and nevertheless the stable
representation and the proved bounds are valid.
\begin{table}[htbp]
\centering\small
\begin{tabular}{rrlll}
\toprule
$q$ & $t$ & Scaled tail $E$ & $I+g/(2t)$ & Error bound \\
\midrule
0 & 20 & 0.969362726 & 0.968260538 & 0.00165490768 \\
0 & 80 & 0.936111703 & 0.936042753 & 0.00010343173 \\
1 & 20 & 0.536929936 & 0.536260165 & 0.00100571307 \\
1 & 80 & 0.517552067 & 0.517510165 & 6.2857067e-5 \\
3 & 20 & 0.048558868 & 0.0484939007 & 9.75587143e-5 \\
3 & 80 & 0.0467372268 & 0.0467331622 & 6.09741964e-6 \\
5 & 20 & 0.00167737852 & 0.00167506571 & 3.47313388e-6 \\
5 & 80 & 0.00161339799 & 0.00161325328 & 2.17070868e-7 \\
\bottomrule
\end{tabular}

\caption{Crossover values and the upper bound in \eqref{eq:trapezoid}.
The approximation is always below the scaled tail. Values are rounded;
the inequalities are theorems, while these evaluations use ordinary
high-precision arithmetic, not directed-rounding intervals.}
\label{tab:crossover}
\end{table}

\section{Reproducibility, proof audit, and formalization targets}
\label{sec:verification}

\subsection{What was actually checked}
The accompanying program \texttt{verification/verify.py} uses exact
rational arithmetic for finite exponent supports and divided-power
polynomials. Its recorded run, with seed $20260929$, reports:
\begin{center}
\begin{tabular}{@{}lr@{}}
\toprule
Check & Count\\
\midrule
Scalar polynomial-inverse identities &105\\
Path-kernel degree and leading coefficient identities &210\\
Finite triangular systems &30\\
Differential-equation entry checks &540\\
Normalized-gauge entry checks &540\\
Constant-connection entry checks &540\\
Weighted path norm bounds &349\\
Explicit resonant-chain formulas &49\\
Crossover parameter cases &96\\
\bottomrule
\end{tabular}
\end{center}

For each finite system, the code independently computes $H$ by path
expansion and by triangular recursion and checks equality. It verifies
\eqref{eq:H-equation}, constructs $(P,\mathcal R)$, verifies
\eqref{eq:P-equation}, and checks that the connection between the two
fundamental matrices is constant after the diagonal monomial shifts.
The spectra include repeated eigenvalues, and the exponent samples include
exact resonances. Finite tests also check the weighted inequality and the
arbitrary-rank chain formula.

The numerical cases use $q\in\{0,1,3,5\}$,
$p\in\{1.5,2,2.5,6.5\}$, $c\in\{1/2,1,2\}$, and $t\in\{20,80\}$.
They evaluate the positive series with analytic series-tail estimates and
check both inequalities in Theorem~\ref{thm:crossover}. The executed
environment was Python 3.13.5 and mpmath 1.3.0 at 75 decimal digits.
Floating-point evaluations are not represented as interval certificates.

These tests do not prove the infinite-support theorem, the closed-graph
argument, or the universal quantifiers. Those are conventional mathematical
proofs in the preceding sections. No Lean theorem implementing this
article's new results was produced or kernel-checked.

\subsection{Logical dependencies and delicate points}
The core proof dependencies are
\[
\begin{aligned}
 \text{finite path fibers}
 &\Longrightarrow\text{formal path formula}\\
 &\Longrightarrow\text{exact weighted criterion}\\
 &\Longrightarrow\text{path gap equivalence}.
\end{aligned}
\]
The normalized log-free conclusion adds only the triangular normal-form
construction and the constant-field comparison. The chain formulas and
crossover are explicit calculations that can be checked independently of
the functional-analytic necessity proof.

Several restrictions are indispensable to the stated argument. Acyclicity
bounds the number of operators in a path and makes the set of paths finite.
Independence of edge inputs permits isolation of one path. Positive real
exponents make the bounded-total small-divisor argument transparent and
fix the near-identity orientation. Well ordering, rather than local
finiteness, provides finite exact coefficient fibers. The target of the
closed-graph argument is a fixed-support Banach space. Exact resonances
are handled by $T_0$ rather than by a fictitious division by zero.

The normalization comparison is an equivalence of membership, not an
identity of operator norms. The all-radii result is stated for the
log-free gauge and for the coefficient norm of its polynomial-logarithmic
counterpart. The basepoint realization is a grouped analytic expression;
its divergent pieces are never summed separately. Finally, the
crossover tail is defined before taking a limit and is not conflated with
a raw Hahn coefficient tail.

\subsection{A realistic Lean decomposition}
The existing support bridge \cite{repo-lean} supplies the two-factor
well-ordering and finite-decomposition interface. A formalization of the
present results would add the following layers.

First, finite acyclic paths and their exponent tuple fibers can be handled
without analysis. The polynomial operators $T_\beta$, their right-inverse
identities, and the divided-power leading-coefficient lemma are
finite-dimensional algebra. These are especially suitable initial targets.

Second, the formal path expansion and the triangular log-free normal form
need explicit Hahn coefficient operations, a strong Euler derivation, and
the polynomial-log constant-field lemma. Their proofs should expose the
support bounds rather than rely on an informal field name.

Third, the weighted theorem requires fixed-support $\ell^1$ spaces,
coordinate continuity, closed graphs, and uniform boundedness. A formal
statement should quantify over independent edge families. The exact
operator-norm equality is a useful interface between the finite path
algebra and the analytic layer.

Finally, the sharp converse needs the bounded-tuple small-divisor sequence
and the leading-coefficient lower bound. The crossover proof can be
formalized separately as a positive-series and trapezoidal-error theorem.
This decomposition distinguishes existing repository support lemmas from
new proof obligations; the former do not automatically verify the latter.

\section{Further research questions}
\label{sec:future}

\subsection{Cycles and infinite matrix-word languages}
For a nontriangular perturbation, matrix words may have arbitrarily large
length. Positive input exponents still bound the length of words below a
fixed spectral cut, but they do not reduce the entire solution to finitely
many multilinear maps. Can a finite critical-word condition, combined with
a large-exponent majorant, characterize universal convergence? The first
nontrivial case is a two-vertex directed cycle with two accumulating
supports. A proof must address repeated occurrences of the same input,
not merely replace ``path'' by ``walk'' in Theorem~\ref{thm:weighted}.

\subsection{Fixed Jordan blocks and zero-action edges}
A fixed nilpotent part in $A_0$ behaves like a collection of zero-exponent
edges. For a triangular constant matrix, finite path length survives, but
some edge coefficients are fixed rather than independently variable.
Determine the exact accessible small-divisor powers in terms of the
Jordan filtration and the variable-edge pattern. The chain in
Section~\ref{sec:chain} suggests that logarithmic height, not merely the
number of near-resonances, is the relevant weight exponent.

\subsection{Coupled matrix directions and relations among words}
The four-dimensional cancellation example shows that a graph alone is
insufficient when coefficients multiply prescribed linear combinations of
matrix units. A natural replacement is a word algebra modulo the matrix
relations, with a kernel attached to each surviving coefficient tensor.
Find an exact independent-parameter criterion that remains valid after
these relations are imposed. Even two fixed nilpotent generators with a
nontrivial quadratic relation would be a useful first classification.

\subsection{Optimal factorized weights}
The exact criterion asks for
$|Q_p(\boldsymbol\alpha)|\leq C_p\prod w_e(\alpha_e)$.
Which minimal or Pareto-optimal families of edge weights satisfy all these
inequalities? For finitely sampled supports, taking logarithms leads to
finite inequality systems; for accumulating supports, one must control
asymptotic factorization of singular path kernels. Such a result would turn
the criterion into a systematic regularity-design tool rather than a
post hoc convergence test.

\subsection{Moving spectra through an accumulation point}
Let a spectral difference move by a parameter $\delta$ through a left
accumulation point. Canonical normalizers have denominators such as
$\delta-1/n$, while basepoint integrals can remain regular in $\delta$.
Determine a uniform theorem when $t/n$ and $\delta t$ compete. Exact
resonance, near resonance, and the accumulation boundary should be treated
in a single analytic representation, not by separately divergent formulas.

\subsection{Multiplicative realization and iterated integrals}
The basepoint construction is compatible with ordinary analytic matrix
products because it solves a genuine differential system. This does not
yet define a product-preserving summation map on the underlying formal
Hahn algebra. Can one identify a restricted algebra of accumulating-support
symbols whose realizations are closed under products and satisfy the
shuffle identities of iterated integrals? The first target should include
two interacting resonant chains, where nonzero products force compatibility
conditions beyond a scalar subtraction rule.

\subsection{Triangular nonlinear systems and tree kernels}
For polynomial triangular differential systems, substitutions branch into
rooted trees rather than simple paths. Does a finite-tree analogue of the
weighted kernel criterion hold, and when do diagonal coefficient
identifications obstruct its necessity? Polarization may recover some
independence, but repeated use of one input and coinciding exponents must
be handled explicitly. This would connect the present classification to
nonlinear transseries reversion without presupposing convergence.

\subsection{Effective presentation of infinite supports}
The criterion is semantic: it refers to left accumulation of exact
sumsets. For which computable descriptions of $E_e$ can the condition be
decided or certified? Finite rational supports are elementary, but arbitrary
computable increasing sequences can encode difficult limiting questions.
A useful next result would distinguish finite symbolic certificates,
semidecidable obstruction searches, and genuinely decidable presentation
classes.

\subsection{Irregular exponential sectors}
Our systems are regular singular in the Euler sense and the expansions
use real powers and logarithms. Irregular systems introduce exponential
sectors and potentially factorial coefficient growth. Can one separate
an accessible finite-spectral accumulation obstruction from a Borel-plane
summability obstruction in one explicit triangular irregular model?
A successful theorem should retain two distinct norms, rather than call
both mechanisms ``resurgence.''

\subsection{Kernel-checked necessity, not only finite examples}
The most valuable formal target is the full implication from universal
convergence to bounded path kernels. Its proof exposes the quantifier that
finite numerical experiments cannot establish. A staged implementation
should first verify the scalar and finite-path algebra, then the
fixed-support closed-graph lemma, and finally the exact gap equivalence.
The supplied rational verifier is a reproducibility aid, not a substitute
for this goal.

\section{Conclusion}

For acyclic matrix patterns, universal absolute convergence is governed by
what the differential system can actually multiply. The complete exponent
semigroup is too coarse, and inspecting each edge separately is too weak.
The correct invariant is the family of endpoint-labelled path sumsets.
The exact weighted refinement retains the full polynomial-logarithmic
kernel and therefore detects powers of small denominators caused by
resonant descendants.

This resolves a concrete fixed-support question from the repository for a
substantial precisely specified class, gives explicit counterexamples to
two natural simplifications, and quantifies an arbitrary-rank resonance
loss. The stable chain realization also demonstrates that divergence of
a canonical normalizer is not nonexistence of analytic solutions or loss
of every finite asymptotic prefix. The remaining questions concern the
extra algebra introduced by cycles, coupled directions, nonlinear trees,
and irregular exponential scales. Those are extensions of the proved
boundary, not assumptions silently included in it.

\clearpage
\begin{thebibliography}{99}

\bibitem{repo-index}
Vladimir Reshetnikov, \emph{ProveIt: Series and transseries}, repository
index at snapshot \snap, inspected 29 September 2026.
\path{Analysis/Transseries/docs/series-and-transseries/README.md}.
\href{https://github.com/VladimirReshetnikov/ProveIt/blob/9250bbf8af80dacf7b252d9d0320b122c80961ba/Analysis/Transseries/docs/series-and-transseries/README.md}{Pinned repository source}.

\bibitem{repo-hahn}
\emph{Finite Resonance Certificates and Sharp Analytic Normalization for
Hahn--Fuchsian Systems}, research article and editorial amendments in
ProveIt, 29 September 2026; particularly Sections 3 and 11.
\path{Analysis/Transseries/docs/series-and-transseries/Hahn_Fuchsian_Resonance_Analytic_Normalization/hahn_fuchsian.tex}.
Snapshot \snap.
\href{https://github.com/VladimirReshetnikov/ProveIt/blob/9250bbf8af80dacf7b252d9d0320b122c80961ba/Analysis/Transseries/docs/series-and-transseries/Hahn_Fuchsian_Resonance_Analytic_Normalization/hahn_fuchsian.tex}{Pinned article source}.

\bibitem{repo-lean}
Vladimir Reshetnikov, \emph{Well-based supports: Dickson's lemma and
Neumann's lemma}, Lean source in ProveIt.
\path{Analysis/Transseries/Lean/Transseries/TransseriesWellBased.lean}.
Snapshot \snap.
\href{https://github.com/VladimirReshetnikov/ProveIt/blob/9250bbf8af80dacf7b252d9d0320b122c80961ba/Analysis/Transseries/Lean/Transseries/TransseriesWellBased.lean}{Pinned Lean source}.

\bibitem{repo-block}
\emph{Resonance-Block Summation of Transseries: Small divisors, coalescing
Stokes actions, and nonlinear inversion}, package README in ProveIt,
29 September 2026.
\path{Analysis/Transseries/docs/series-and-transseries/Resonance_Block_Summation_Transseries/README.md}.
Snapshot \snap.
\href{https://github.com/VladimirReshetnikov/ProveIt/blob/9250bbf8af80dacf7b252d9d0320b122c80961ba/Analysis/Transseries/docs/series-and-transseries/Resonance_Block_Summation_Transseries/README.md}{Pinned package documentation}.

\bibitem{vdh}
Joris van der Hoeven, \emph{Transseries and Real Differential Algebra},
Lecture Notes in Mathematics 1888, Springer, 2006.
\href{https://www.texmacs.org/joris/ln/ln.html}{Author-hosted text}.

\bibitem{iy}
Yulij Ilyashenko and Sergei Yakovenko,
\emph{Lectures on Analytic Differential Equations}, Graduate Studies in
Mathematics 86, American Mathematical Society, 2008.
\href{https://www.wisdom.weizmann.ac.il/~yakov/book-titlepage.htm}{Author-hosted bibliographic information}.

\bibitem{yakovenko}
Sergei Yakovenko, \emph{Local theory of Fuchsian systems}, Lecture 3,
10 November 2014. Author's course notes on resonant normal forms and their
fundamental matrices.
\href{https://yakovenko.wordpress.com/2014/11/10/lecture-3-nov-10-2014/}{Author's lecture page}.

\bibitem{gg}
R. R. Gontsov and I. V. Goryuchkina,
\emph{On the convergence of generalized power series satisfying an
algebraic ODE}, Asymptotic Analysis 93 (2015), 311--325.
\href{https://doi.org/10.3233/ASY-151297}{doi:10.3233/ASY-151297};
\href{https://arxiv.org/abs/1407.1330}{arXiv:1407.1330}.

\bibitem{ggl}
Renat Gontsov, Irina Goryuchkina, and Alberto Lastra,
\emph{Small divisors in the problem of the convergence of generalized
power series solutions of $q$-difference equations}, 2022 preprint.
\href{https://arxiv.org/abs/2209.09365}{arXiv:2209.09365}.

\end{thebibliography}
\end{document}
